\documentclass[10pt]{article}
\usepackage[margin=0.86in]{geometry}
\usepackage[T1]{fontenc}
\usepackage{lmodern,microtype}
\usepackage{amsmath,amssymb}
\usepackage[colorlinks=true,allcolors=blue!55!black]{hyperref}
\setlength{\parindent}{0pt}
\setlength{\parskip}{0.55em}
\begin{document}

\begin{center}
{\Large\bfseries Literature answer to the nonsingular quartic flat-extension question}\par
\medskip
{\small arXiv:0908.3230, Question 3.5 $\longrightarrow$ arXiv:1412.7882, Theorem 2.1}
\end{center}

\textbf{Status: literature\_already\_answered.}

\textbf{Original source.} Lawrence Fialkow and Jiawang Nie,
\emph{Positivity of Riesz Functionals and Solutions of Quadratic and Quartic Moment Problems}, arXiv:0908.3230.  Question 3.5 on page 12 of the arXiv PDF asks, for a bivariate quartic sequence $y\in\mathcal M_{2,4}$ with $M_2(y)\succ0$:
\begin{quote}
Does $M_2(y)$ have a flat extension?  Does $y$ have a degree-six extension $\widetilde y$ such that $M_3(\widetilde y)$ is positive and has a flat extension?
\end{quote}

\textbf{Decisive later answer.} Ra\'ul E. Curto and Seonguk Yoo,
\emph{Concrete Solution to the Nonsingular Quartic Binary Moment Problem}, arXiv:1412.7882; Proc. Amer. Math. Soc. \textbf{144} (2016), 249--258.  Theorem 2.1 on page 3 states:
\begin{quote}
If $M(2)$ is positive and invertible, then it has a representing measure with exactly six atoms; equivalently, $M(2)$ admits a flat extension $M(3)$.
\end{quote}

This is exactly the positive-definite $6\times6$ matrix in Question 3.5.  A six-atomic representing measure makes the degree-three moment matrix have rank at most six, while its $M_2$ principal block already has rank six; hence the extension is flat.  Curto--Yoo explicitly cite the Fialkow--Nie paper and contrast their constructive six-atom theorem with its earlier existential bound of fifteen atoms.  The later paper therefore knowingly treats the same nonsingular bivariate quartic problem, rather than merely implying an unrelated special case.

\textbf{Scope.} Both parts of Question 3.5 are answered affirmatively.  This status does not claim to answer Fialkow--Nie's much broader Question 1.2 for arbitrary degree, dimension, and support set.

\textbf{Evidence and verification.}\par
Searches by the exact question, title, ``flat extension'', and ``six-atomic'' located arXiv:1412.7882.  The source Question 3.5 and supporting Theorem 2.1 were checked in the official PDFs saved with this note.  The final one-page PDF was compiled and visually inspected.

\end{document}
